\documentclass[11pt]{article}
\usepackage[utf8]{inputenc}
\usepackage{amsmath,amssymb}
\usepackage{booktabs}
\usepackage{graphicx}
\usepackage{multirow}
\usepackage[margin=1in]{geometry}
\usepackage{hyperref}
\usepackage{xcolor}

\title{Exception Events Are All You Need: Fault Indicator Isolation\\for Verdict Prediction in Code Execution Traces}

\author{Anonymous}
\date{}

\begin{document}
\maketitle

\begin{abstract}
Execution traces contain rich runtime information, yet their effectiveness for verdict prediction (pass/fail classification of code solutions) remains inconsistent across models and representations. We systematically ablate trace components across four model configurations (Qwen-2.5-Coder-7B, Qwen-2.5-Coder-1.5B, CodeLlama-13B-Instruct, DeepSeek-Coder-6.7B) and discover that \textbf{exception events}---compact structured error summaries---are sufficient for near-perfect verdict prediction, achieving 99.33--99.83\% accuracy across all models with cross-model standard deviation of only 0.24 percentage points. In contrast, VOTA traces that incorporate behavioral patterns (loop compression) show high model-dependency (cross-model std = 6.56pp, range 84.45--99.41\%). Component ablation reveals a strict two-tier diagnostic hierarchy: Tier~1 signals (exceptions) are individually near-perfect and model-invariant, while Tier~2 signals (loop patterns) are useful but seed-sensitive (std = 3.47pp across seeds). Removing structural formatting markers ([ERROR] tags) causes only 0.50$\pm$0.17pp degradation, confirming that exception \emph{content} carries the fault signal. A label-only loss ablation decomposes the full-trace-to-VOTA gap: 26\% gradient dilution, 74\% information overload. Our results establish that verdict prediction is fundamentally a fault detection task, and the design principle for verdict-optimized traces is diagnostic signal isolation, not generic compression.
\end{abstract}

\section{Introduction}

Execution traces---runtime logs capturing variable states, control flow, and exception events---are a natural source of supervision for code understanding tasks. Recent work has explored trace-augmented training for code generation (DCSH, EMNLP Findings 2025; TracePile, NeurIPS 2025), finding inconsistent benefits across representations. TRACED (ICSE 2024) demonstrated that traces help encoder-based test verdict classification, but did not identify which trace components drive performance.

We address a focused question: \emph{what is the minimal sufficient diagnostic signal for verdict prediction?} Through systematic component ablation on a dataset of 5,910 code solutions (591 problems $\times$ 10 LLM-generated solutions from HumanEval + MBPP-sanitized), we make three contributions:

\begin{enumerate}
\item \textbf{Exception sufficiency}: Exception events alone match or exceed full VOTA traces, achieving 99.83\% on Qwen-7B (3 seeds, $\sigma$=0.00) and $\geq$99.33\% on three additional model configurations---with $\sim$16$\times$ fewer tokens and no execution reasoning requirement.
\item \textbf{Information density bottleneck}: Label-only loss masking quantifies that 74\% of the full-trace disadvantage stems from information overload (not gradient dilution), explaining why compression helps.
\item \textbf{Two-tier diagnostic hierarchy}: Tier~1 (exception events) is model-invariant and seed-robust; Tier~2 (loop behavioral patterns) is useful but exhibits seed-dependent stability (std=3.47pp), fully subsumed by Tier~1.
\end{enumerate}

\section{Method}

\subsection{Task and Dataset}

Given a programming problem, a candidate solution, and (optionally) an execution trace representation, the model predicts a binary verdict: \texttt{pass} or \texttt{fail}. We construct our dataset from HumanEval (164 problems) and MBPP-sanitized (427 problems), generating 10 LLM solutions per problem (5,910 total). We split by problem into train (4,720 samples) and test (1,190 samples) with zero input overlap.

\subsection{Trace Representations}

We evaluate five trace conditions of decreasing information density:

\begin{itemize}
\item \textbf{no\_trace}: Source code + problem description only (baseline).
\item \textbf{full\_trace}: Complete execution trace with line-level variable states, control flow, and exceptions.
\item \textbf{VOTA}: Verdict-Optimized Trace Abstraction---a compressed trace retaining three components: (1) exception summaries with \texttt{[ERROR]} markers, (2) actual/expected value pairs, (3) loop iteration compression.
\item \textbf{exception\_only}: Only the exception summary component (\texttt{[ERROR] ExceptionType: message}).
\item \textbf{exception\_only\_nomarker}: Exception content (type + message) without the \texttt{[ERROR]} structural marker.
\end{itemize}

Additionally, we evaluate isolated components: \textbf{loop\_only} (only loop compression patterns).

\subsection{Training Setup}

All models are fine-tuned with LoRA (rank=16, $\alpha$=32) for 3 epochs on the SFT objective. We use max sequence length 2048 tokens. Each condition is evaluated across 3 random seeds (42, 123, 456) to assess training stability. Models: Qwen-2.5-Coder-7B (primary), Qwen-2.5-Coder-1.5B (scale ablation), CodeLlama-13B-Instruct (architecture diversity), DeepSeek-Coder-6.7B (from prior work).

\section{Experiments}

\subsection{Experiment 1: Nomarker Ablation---Is Structural Formatting Required?}

We test whether the \texttt{[ERROR]} marker is an active ingredient or whether exception \emph{content} alone suffices.

\begin{table}[h]
\centering
\caption{Effect of removing [ERROR] format marker (Qwen-7B, 3 seeds).}
\label{tab:nomarker}
\begin{tabular}{lcccc}
\toprule
Seed & Accuracy & Errors & Drop vs.\ marker \\
\midrule
42 & 99.33\% & 8 & $-$0.50pp \\
123 & 99.16\% & 10 & $-$0.67pp \\
456 & 99.50\% & 6 & $-$0.33pp \\
\midrule
mean$\pm$std & 99.33\%$\pm$0.17 & 8.0$\pm$2.0 & $-$0.50$\pm$0.17pp \\
\bottomrule
\end{tabular}
\end{table}

The marker condition achieves 99.83\% (2 errors, all seeds identical). Removing the \texttt{[ERROR]} marker introduces only 0.50$\pm$0.17pp degradation (Table~\ref{tab:nomarker}). This confirms that exception \emph{content itself}---the exception type and message---carries the fault signal. The structural formatting marker provides at most a marginal alignment benefit, not a necessary cue. A regex baseline detecting \texttt{[ERROR]} presence achieves only 93.45\%, far below SFT performance, ruling out trivial keyword matching.

\subsection{Experiment 2: Loop-Only Component---Seed-Dependent Stability}

We isolate the Tier~2 behavioral signal (loop iteration patterns) to characterize its reliability.

\begin{table}[h]
\centering
\caption{Loop-only component across 3 seeds (Qwen-7B).}
\label{tab:loop}
\begin{tabular}{lccl}
\toprule
Seed & Accuracy & Fail Recall & Interpretation \\
\midrule
42 & 85.63\% & 0.820 & baseline \\
123 & 81.18\% & 0.748 & worse LoRA basin \\
456 & 88.07\% & 0.851 & better LoRA basin \\
\midrule
mean$\pm$std & 84.96\%$\pm$3.47 & 0.806$\pm$0.053 & seed-sensitive \\
\bottomrule
\end{tabular}
\end{table}

Loop-only achieves mean 84.96\% (Table~\ref{tab:loop})---substantially above no\_trace (75.13\%) but far below exception\_only (99.83\%). Critically, the cross-seed standard deviation is 3.47pp (range: 81.18--88.07\%), compared to 0.00pp for exception\_only. The loop signal's learnability is LoRA-initialization-dependent: seed~123 falls into a suboptimal basin where fail recall drops to 0.748 (vs.\ 0.851 for seed~456). This seed sensitivity is a defining characteristic of Tier~2 signals---useful but unreliable without training stabilization.

The loop-only mean (84.96\%) approximates label-only full\_trace (86.22\%), confirming that VOTA's loop compression precisely extracts the useful behavioral subset from full traces.

\subsection{Experiment 3: Cross-Scale Robustness (Qwen 1.5B)}

We test whether exception signal effectiveness depends on model scale.

\begin{table}[h]
\centering
\caption{Cross-scale comparison: Qwen-2.5-Coder 1.5B vs 7B.}
\label{tab:scale}
\begin{tabular}{lccr}
\toprule
Condition & 1.5B & 7B & Gap \\
\midrule
no\_trace & 67.31\% & 75.13\% & $-$7.82pp \\
full\_trace & 74.45\% & 81.76\% & $-$7.31pp \\
VOTA & 97.00\% & 99.41\% & $-$2.41pp \\
exception\_only & 99.83\% & 99.83\% & 0.00pp \\
\bottomrule
\end{tabular}
\end{table}

Exception\_only achieves \emph{identical} performance (99.83\%) regardless of model scale (Table~\ref{tab:scale}). In contrast, VOTA degrades by 2.41pp on the smaller model, and full\_trace/no\_trace degrade by $\sim$7pp. The cross-scale gap narrows monotonically with information density: as the signal becomes more targeted, model capacity matters less. At the extreme (exception\_only), even 1.5B parameters suffice for near-perfect fault detection.

The 1.5B full\_trace result (74.45\%) shows that information overload affects smaller models similarly to larger ones (+7.1pp lift vs.\ no\_trace, comparable to 7B's +6.6pp), but the full-trace-to-exception lift is \emph{larger} for 1.5B (25.4pp vs.\ 18.1pp for 7B)---smaller models benefit more from signal extraction.

\subsection{Experiment 4: Cross-Architecture Validation (CodeLlama-13B)}

We validate on a third architecture (CodeLlama-13B-Instruct) with different tokenizer, pretraining data, and model family.

\begin{table}[h]
\centering
\caption{CodeLlama-13B-Instruct results.}
\label{tab:codellama}
\begin{tabular}{lccc}
\toprule
Condition & Accuracy & Pass F1 & Fail F1 \\
\midrule
no\_trace & 76.22\% & 0.826 & 0.631 \\
VOTA & 92.44\% & 0.951 & 0.849 \\
exception\_only & 99.50\% & 0.996 & 0.992 \\
\bottomrule
\end{tabular}
\end{table}

Exception\_only achieves 99.50\% on CodeLlama (Table~\ref{tab:codellama}), confirming cross-architecture robustness. The exception-to-VOTA gap is 7.06pp---intermediate between Qwen-7B (0.42pp) and DeepSeek-6.7B (14.88pp), consistent with CodeLlama's intermediate execution reasoning capability (estimated CRUXEval-I $\sim$46\%).

\subsection{Experiment 5: Full Cross-Model Summary}

\begin{table}[h]
\centering
\caption{Cross-model comparison across all architectures. $\dagger$DeepSeek median of 3 seeds (seed42/456=99.33\%, seed123=91.09\%).}
\label{tab:crossmodel}
\begin{tabular}{lccccc}
\toprule
Condition & Qwen-7B & Qwen-1.5B & CodeLlama-13B & DeepSeek-6.7B & Std \\
\midrule
exception\_only & 99.83\% & 99.83\% & 99.50\% & 99.33\%$\dagger$ & 0.24pp \\
VOTA & 99.41\% & 97.00\% & 92.44\% & 84.45\% & 6.56pp \\
no\_trace & 75.13\% & 67.31\% & 76.22\% & 73.36\% & 3.86pp \\
\bottomrule
\end{tabular}
\end{table}

Table~\ref{tab:crossmodel} presents the central result. Exception\_only achieves cross-model std of only 0.24pp (range 99.33--99.83\%), while VOTA std is 6.56pp (range 84.45--99.41\%)---a 27$\times$ ratio. Exception signal effectiveness is nearly model-invariant; VOTA effectiveness is highly model-dependent. All four model configurations show exception\_only $\geq$ VOTA.

The VOTA cross-model gap (14.96pp between Qwen-7B and DeepSeek) has two contributing factors: (1) model execution reasoning capability differences (CRUXEval-I: Qwen 56.5\% vs DeepSeek 39.0\%), and (2) training-time trace truncation on DeepSeek (91.7\% of VOTA samples truncated from mean 2626 to $\sim$2048 tokens; labels preserved but trace content lost). Group analysis confirms truncation causality: train-truncated samples achieve 83.04\% vs.\ train-intact samples at 100\% (Fisher exact $p \approx 0$).

\subsection{Experiment 6: VOTA vs Exception Gap by Model Capability}

\begin{table}[h]
\centering
\caption{VOTA-to-exception gap correlates with model execution reasoning capability.}
\label{tab:gap}
\begin{tabular}{lccc}
\toprule
Model & exception\_only & VOTA & Gap \\
\midrule
Qwen-2.5-Coder-7B & 99.83\% & 99.41\% & 0.42pp \\
CodeLlama-13B & 99.50\% & 92.44\% & 7.06pp \\
DeepSeek-Coder-6.7B & 99.33\% & 84.45\% & 14.88pp \\
\bottomrule
\end{tabular}
\end{table}

The gap between VOTA and exception\_only widens as model execution reasoning decreases (Table~\ref{tab:gap}). Models with weaker trace reasoning struggle to extract useful signal from VOTA's Tier~2 components, while exception signal---requiring no execution reasoning---remains uniformly effective. This validates the design principle: simpler signals are more robust across the model capability spectrum.

\section{Analysis}

\paragraph{Information Density Decomposition.}
On Qwen-7B, full\_trace achieves 81.76\% while VOTA achieves 99.41\% (gap = 17.65pp). Label-only loss masking (which eliminates gradient dilution by computing loss only on verdict tokens) achieves 86.22\%---closing 4.5pp of the gap (26\%). The remaining 13.0pp (74\%) is attributable to information overload: even with oracle loss, the model cannot extract the relevant signal from noisy full traces. Only trace compression (information density improvement) resolves this bottleneck.

\paragraph{Irreducible Error.}
On Qwen-7B exception\_only, all 3 seeds produce identical errors on exactly 2 samples (ids 901, 905), both from HumanEval/95 (\texttt{check\_dict\_case}). Root cause: Python's \texttt{all()} short-circuits on a generator, triggering \texttt{GeneratorExit} (a \texttt{BaseException} subclass, not a real error). The tracer marks this as \texttt{[ERROR]}, and the model reasonably but incorrectly predicts fail. These 2 samples are the only \texttt{GeneratorExit} instances in the test set, with no training signal to learn their benign nature. Thus 99.83\% represents the theoretical ceiling for this representation.

\paragraph{Why Exception Content Suffices Without Markers.}
The nomarker ablation (Section~3.1) shows only 0.50pp drop. Combined with the regex baseline's 93.45\% (vs.\ SFT's 99.83\%), this demonstrates that the model learns from exception \emph{semantics}---the type name and message describe the fault's nature---rather than relying on \texttt{[ERROR]} as a keyword trigger. The formatting marker provides a mild alignment benefit (easier pattern recognition during early training) but is not the active ingredient.

\section{Related Work}

\textbf{Trace-augmented code training.} DCSH (Chen et al., EMNLP Findings 2025) evaluated five trace formats for code generation, finding inconsistent benefits. TracePile (Ding et al., NeurIPS 2025) scaled trace pretraining to 2.6M samples but evaluated only generation tasks. Neither work studied verdict prediction or identified which trace components drive classification.

\textbf{Test verdict prediction.} TRACED (Luo et al., ICSE 2024) showed execution traces help encoder-based verdict classification but did not perform component ablation. Our work complements TRACED by identifying exception events as the minimal sufficient signal and demonstrating this across decoder-based models via SFT.

\textbf{Fault localization.} Classical fault localization uses execution spectra and statistical methods. Our finding that exception events suffice connects to the fault localization intuition that crash points are the strongest diagnostic signal, now validated in the neural SFT setting.

\section{Conclusion}

We present a systematic study of trace components for verdict prediction, establishing that exception events are the minimal sufficient diagnostic signal. Key findings: (1) Exception\_only achieves 99.33--99.83\% across four model configurations with 0.24pp cross-model std, matching or exceeding VOTA with $\sim$16$\times$ fewer tokens and no execution reasoning requirement. (2) Information overload accounts for 74\% of the full-trace disadvantage; gradient dilution is secondary. (3) A two-tier hierarchy exists: Tier~1 (exceptions) is model-invariant and seed-robust; Tier~2 (loop patterns) shows seed-dependent stability (std=3.47pp). These results reframe VOTA's value as fault indicator isolation rather than generic compression, and establish a clear design principle for verdict-optimized traces: extract the highest-specificity diagnostic moment; discard everything else.

\paragraph{Limitations.} Our evaluation covers a single task (verdict prediction) on one dataset (HumanEval + MBPP-sanitized). While four model architectures provide diversity, broader task and dataset validation is needed. DeepSeek exhibits seed-dependent absolute performance on exception\_only (2/3 seeds at 99.33\%, 1 outlier at 91.09\%), suggesting that training stabilization remains important for some models.

\end{document}
